\documentclass[10pt,a4paper]{article}

\usepackage[english]{babel}
\usepackage[utf8]{inputenc}
\usepackage[T1]{fontenc}
\usepackage[a4paper, margin=2cm]{geometry}
\usepackage{amsmath,amsfonts,amssymb}
\usepackage[colorlinks]{hyperref}
\usepackage{multicol}
\usepackage{graphicx}
\usepackage{float}
\usepackage{booktabs}

\usepackage[round]{natbib}
\bibliographystyle{abbrvnat}

\usepackage{fancyhdr}
\pagestyle{fancy}
\fancyhead[L]{02460-2026 Mini-project 2}
\fancyhead[R]{s039894, s271828 \& s314159}

% For including code
% From https://www.overleaf.com/learn/latex/Code_listing
\usepackage{listings}
\usepackage{xcolor}

\definecolor{codegreen}{rgb}{0,0.6,0}
\definecolor{codegray}{rgb}{0.5,0.5,0.5}
\definecolor{codepurple}{rgb}{0.58,0,0.82}
\definecolor{backcolour}{rgb}{0.95,0.95,0.92}

\lstdefinestyle{mystyle}{
    backgroundcolor=\color{backcolour},   
    commentstyle=\color{codegreen},
    keywordstyle=\color{magenta},
    numberstyle=\tiny\color{codegray},
    stringstyle=\color{codepurple},
    basicstyle=\ttfamily\footnotesize,
    breakatwhitespace=false,         
    breaklines=true,                 
    captionpos=b,                    
    keepspaces=true,                 
    numbers=left,                    
    numbersep=5pt,                  
    showspaces=false,                
    showstringspaces=false,
    showtabs=false,                  
    tabsize=2,
    upquote=true
}

\lstset{language=Python, style=mystyle}

\usepackage{lipsum}

\begin{document}

%%% Main text part %%%

\section*{Mini-project 2 in Advanced Machine Learning (02460)}
by Alice (s039894), Bob (s271828) \& Carlo (s314159) on 3 April 2025

\begin{multicols}{2}
    \subsection*{Part A: pullback geodesics}
    \lipsum[1] We show some plots in figure~\ref{fig:geodesics}.
%    \begin{figure}[H]
%        \centering
%        \includegraphics[width=0.49\linewidth]{geodesics.pdf}
%        \includegraphics[width=0.49\linewidth]{geodesics.pdf}
%        \caption{Geodesics between randomly sampled latent variables. The left panel shows results using a standard VAE, while the right panel shows an ensemble decoder VAE with 3 ensemble members.}
%        \label{fig:geodesics}
%    \end{figure}
    \lipsum[2]

    \subsection*{Part B: Ensemble VAE geometry}
    \lipsum[1]
%        \begin{figure}[H]
%        \centering
%        \includegraphics[width=0.6\linewidth]{cov.pdf}
%        \caption{The coefficient of variation across multiple training runs for different number of decoders.}
%        \label{fig:plot}
%    \end{figure}

    \lipsum[1]
\end{multicols}

\newpage

%%% Code snippet part %%%

\section*{Code snippets}
The following beautiful code snippet prints \lstinline{"Hello world"} 10 times:
\begin{lstlisting}
for i in range(10):
    print("Hello world")
\end{lstlisting}

Here is a more advanced code example that prints \lstinline{"Hello world"} 100 times:
\begin{lstlisting}
for i in range(100):
    print("Hello world")
\end{lstlisting}

\newpage

%%% Bibliography %%%

\bibliography{template}

\end{document}
